\scanpage{51}

% B01 | head, lemma and vertical rule
\uline{Lemme 1} Soit $k$ un corps algébriquement clos, $K$ une extension de type fini, de degré de transcendance $1$, de $k$. Alors $K$ est \deleted{\unsure{algèb|algéb}} quasi-algébriquement clos

% B02 | upper, first corollary
\uline{Corollaire 1} $\widehat H^i(K^{\unsure{*|2|'}},G_m)=0$ pour tout $i$

% B03 | upper, proof
En effet, $K$ et ses extensions galoisiennes $K'$ sont quasi-alg. clos, donc \deleted{\unsure{$H^i(K,G$}} la conclusion est bien connue.

% B04 | upper middle, second corollary
\uline{Corollaire 2} $\deleted{\widehat{\vphantom H}}H^i(K,\Gamma)=0$ \deleted{si $i>0$} \inserted{pour tt $i>1$ et} \deleted{pour} tout groupe abélien fini $\Gamma$ où $G_K$ opère trivialement.

% B05 | middle, roots of unity
Soit en effet $\widetilde K$ \deleted{la} \inserted{une} clôture galoisienne de $K$, [$U\subset\widetilde K^*$ le \deleted{\unsure{sou|un|sous-}} groupe des racines de l'unité dans $\widetilde K$. Alors $\widetilde K^*/U$ est uniquement divisible sauf pour $p$ (\unsure{l'exposant|le apparent} car.) d'où la suite exacte \uline{mod} $p$

% B06 | middle, comparison across the fold
\[\deleted{\unsure{\simeq|=}}\widehat H^i(K,\deleted{\unsure{G_m|\widetilde K}}U)\xrightarrow{\sim}\widehat H^i(K,G_m)\deleted{\unsure{\mathrm{est}|0}}\qquad\text{\uuline{mod $p$}}\]

% B07 | left of comparison, two slightly slanting lines
% slanted
\margin{a-t-on un vrai isomorphisme?}

% B08 | below fold, nullity and its revision
Or le \deleted{premier} deuxième mb est nul, donc le premier l'est, donc comme le premier \deleted{\unsure{sing|groupe}} est $p$-composante nulle, il est nul.] On \deleted{\unsure{prouve|pourra}} \deleted{\unsure{en vertu de|en vertu}} voit plus généralement que si $n$ est premier à la caractéristique, alors

% B09 | lower third, power map
$\widetilde K\xrightarrow{x\to x^n}\widetilde K^{\unsure{n|\times}}$ est un épimorphisme de \deleted{groupe de} noyau $\simeq\mathbb Z/n$, donc la suite exacte de cohomologie donne le résultat $H^i(K,\mathbb Z_n)=0$.

% B10 | foot, deleted heading and characteristic p
\deleted{\uline{Lemme 2} Soit} Reste à prouver $H^i(K,\mathbb Z/p)=0$ si $i>1$. Pour cela on considère la suite exacte $0\longrightarrow\mathbb Z/p\longrightarrow\widetilde K\xrightarrow{\wp}\widetilde K\unsure{)|'}\longrightarrow0$.

% archive sheet number: 50

